\section{第二阶段：VLA 与机器人基座模型}

上一章讲现有基础模型如何进入机器人，本章进入第二阶段：直接为机器人预训练 foundation model。VLA 的目标是把视觉、语言和动作放进一个模型中，让模型从观察和指令直接产生机器人动作。陈建宇用“人就是很智能的 VLA Agent”来解释：人看到世界、理解语言、输出动作，本身就是通用 VLA。

\begin{figure}[H]
\centering
\includegraphics[width=0.92\textwidth]{ch13-02358.jpg}
\caption{VLA Model：将视觉和语言输入映射到 action 输出。投屏帧，约 00:39:18。}
\end{figure}

\begin{knowledgebox}{读图：VLA 的 A 是真正难点}
图中从 vision 和 language 到 action，展示的是 VLA 的最小形态。V 和 L 可以借用 VLM/LLM 的进展，但 A 需要处理频率、连续控制、轨迹、稳定性和本体差异。
\end{knowledgebox}

\subsection{ALOHA 与 Mobile ALOHA：低成本硬件和示范数据}

本节从 VLA 概念回到数据和硬件底座。原因是，机器人基座模型不是只靠论文里的大模型结构就能出现，它首先需要便宜、可复制、能稳定采集示范的系统。ALOHA 在这里承担的是“把真实双手操作数据做出来”的角色，为后续讨论 action chunk、模仿学习和跨任务泛化提供具体参照。

ALOHA 系列不是严格意义上的现代 VLA，但它对机器人学习很重要：低成本双臂硬件、teleoperation 示范、Action Chunking Transformer，以及让机器人模仿细粒度双手操作。Mobile ALOHA 加上移动底盘，把双臂操作扩展到更复杂的移动场景，并以大量演示视频出圈。

\begin{figure}[H]
\centering
\includegraphics[width=0.96\textwidth]{ch15-02658.jpg}
\caption{ALOHA：Action Chunking Transformer 用于双手操作动作序列预测。投屏帧，约 00:44:18。}
\end{figure}

\begin{knowledgebox}{读图：Action Chunking 是动作序列预测}
图中展示 ACT 的结构：输入观测，输出未来一段动作序列。关键是 chunking：模型不是只预测下一步动作，而是预测一段动作轨迹，再滚动执行。这和模型预测控制有相似直觉。
\end{knowledgebox}

\begin{figure}[H]
\centering
\includegraphics[width=0.92\textwidth]{ch16-02943.jpg}
\caption{Mobile ALOHA：将低成本双臂操作扩展到移动全身遥操作。投屏帧，约 00:49:03。}
\end{figure}

\begin{warningbox}{ALOHA 的局限}
ALOHA 很强，但更像高质量模仿学习系统。它仍然偏任务和数据驱动，不等于能跨大量机器人和任务泛化的 VLA 基座模型。
\end{warningbox}

\subsection{Gato、RT-1 与 Octo}

接下来这一组工作把问题从“能否收集高质量示范”推进到“能否训练一个更通用的策略”。从 Gato 到 RT-1，再到 Octo，核心变化是通用代理的想法逐步落到真实机器人数据和开源 policy 框架上。老师强调，这些工作不一定都是今天意义上的 VLA，但它们共同铺出了统一建模、多任务训练和机器人控制规模化的路径。

Gato 是“一个模型处理多种任务”的早期通用代理尝试，思想非常早，但当时模型、数据和算力还不够成熟。RT-1 则更直接面向机器人控制，在真实世界控制数据上训练 Robotics Transformer。Octo 是开源通用机器人策略，强调把不同任务和数据放进一个可复用的 policy 框架。

\begin{figure}[H]
\centering
\includegraphics[width=0.92\textwidth]{ch17-03097.jpg}
\caption{Gato：A Generalist Agent，把多种任务 token 化后统一建模。投屏帧，约 00:51:37。}
\end{figure}

\begin{knowledgebox}{读图：Gato 早在思想上接近 VLA}
Gato 图中显示多任务、多模态输入统一进一个模型。它的历史价值是提出 generalist agent 的方向，但还没有今天 VLM/VLA 的模型规模、数据和机器人执行能力。
\end{knowledgebox}

\begin{figure}[H]
\centering
\includegraphics[width=0.94\textwidth]{ch18-03372.jpg}
\caption{RT-1：Robotics Transformer for real-world control at scale。投屏帧，约 00:56:12。}
\end{figure}

\begin{knowledgebox}{读图：RT-1 是机器人数据规模化的重要节点}
RT-1 把任务说明、图像观测和动作 token 放进 Transformer。它的贡献不只是模型结构，还在于真实机器人控制数据和可扩展训练流程。
\end{knowledgebox}

\begin{figure}[H]
\centering
\includegraphics[width=0.92\textwidth]{ch19-03650.jpg}
\caption{Octo：开源通用机器人策略，支持 task 和 observation token 化。投屏帧，约 01:00:50。}
\end{figure}

\subsection{跨本体学习：CrossFormer 与 GR 系列}

通用机器人模型不能只会一种机械臂或一种任务。CrossFormer 关注 cross-embodied learning，希望一个 policy 覆盖 manipulation、navigation、locomotion 和 aviation 等不同本体。字节 AI Lab 的 GR-1/GR-2 则把视频生成预训练、video-language-action 和 web-scale knowledge 引入机器人操作，强调从大规模视频和知识中学习。

\begin{figure}[H]
\centering
\includegraphics[width=0.94\textwidth]{ch20-03892.jpg}
\caption{CrossFormer：Scaling cross-embodied learning，一种 policy 面向多类本体。投屏帧，约 01:04:52。}
\end{figure}

\begin{knowledgebox}{读图：跨本体是 VLA 泛化的关键}
CrossFormer 图示强调不同机器人本体之间的共享策略。真正通用的机器人基座模型不能只服务单一硬件，否则仍然会退回“一个机器人一个模型”。
\end{knowledgebox}

\begin{figure}[H]
\centering
\includegraphics[width=0.94\textwidth]{ch21-04284.jpg}
\caption{GR-1：基于视频生成预训练的视觉机器人操作工作。投屏帧，约 01:11:24。}
\end{figure}

\subsection{本章小结}

第二阶段开始把机器人数据、动作输出和跨本体泛化放到核心位置。ALOHA 提供低成本示范和动作序列学习，RT-1/Octo/RT-X/OpenVLA 等工作把 VLA 推向更通用的机器人策略。

